% GENERATED FILE. Edit canonical lessons, then run npm run generate:books.
% canonical-source-hash: fnv1a64:394bbed8aec62adf
% canonical-lessons: LA-C454-dextrorsum, LA-C454-sinistrorsum, LA-C454-obliquus, LA-C454-circumeo, LA-C454-circuitus, LA-R454-review-447-454

\chapter{To the right, To the left, Slanting, To go round, Detour}
\label{ch:la-to-the-right-to-the-left-slanting-to-go-round-detour}
\hlchaptermodality{454}

\begin{chapteropening}
\textbf{By the end of this chapter:} \emph{I can give directions: to the right, to the left, at a slant, going round, and taking a detour.}

Complete the last lesson of chapter 454: I can give directions: to the right, to the left, at a slant, going round, and taking a detour.
\end{chapteropening}

\section[dextrōrsum]{dextrōrsum — to the right, towards the right}

\label{lesson:LA-C454-dextrorsum}



\begin{quote}
Before the new one: say the Latin for to find, then the Latin for to track.
\end{quote}

\subsection*{You'll want to know: dextrōrsum}
\textbf{dextrōrsum} — \textquotedblleft{}to the right, towards the right\textquotedblright{}.

\begin{quote}
Post pontem via dextrōrsum it.
\end{quote}

\begin{quote}
\emph{After the bridge the road goes to the right.}
\end{quote}

An adverb, also \emph{dextrōrsus}, from \emph{dexter} 'right' and \emph{vorsum} 'turned'.

\begin{cousinweb}
\emph{dextrōrsum} is a shortened \emph{dextrō-vorsum}, built from \emph{dexter} 'right' and \emph{vorsum}, an old form of \emph{versum} 'turned', from \emph{vertō} 'to turn' (Lewis \& Short).
\end{cousinweb}

\begin{grammarlens}[title={What you've built}]
One more word for telling, explaining and describing at B1.
\end{grammarlens}

\subsection*{Your turn}
\begin{itemize}
\raggedright
  \item \emph{Say it:} \emph{dextrōrsum}
  \item \emph{Say it:} \emph{dextrōrsum}, once more
  \item \emph{Say it:} a sentence of your own with \emph{dextrōrsum}
  \item \emph{Recall:} say the Latin for to find, then the Latin for to track, then say \emph{dextrōrsum} again
\end{itemize}

\begin{tcolorbox}[breakable,colback=teal!4,colframe=teal!35!black,title={Before you move on}]
What does dextrōrsum mean? (\textquotedblleft{}To the right, towards the right\textquotedblright{}.) Say it once more.
\end{tcolorbox}
\section[sinistrōrsum]{sinistrōrsum — to the left, towards the left}

\label{lesson:LA-C454-sinistrorsum}



\begin{quote}
Before the new one: say the Latin for to track, then the Latin for to the right.
\end{quote}

\subsection*{You'll want to know: sinistrōrsum}
\textbf{sinistrōrsum} — \textquotedblleft{}to the left, towards the left\textquotedblright{}.

\begin{quote}
Ad compitum sinistrōrsum flectimus et ad flūmen venīmus.
\end{quote}

\begin{quote}
\emph{At the crossroads we turn to the left and come to the river.}
\end{quote}

An adverb, also \emph{sinistrōrsus}, built like \emph{dextrōrsum}.

\begin{cousinweb}
\emph{sinistrōrsum} is a shortened \emph{sinistrō-vorsum}, built from \emph{sinister} 'left' and \emph{vorsum} 'turned', from \emph{vertō} 'to turn' (Lewis \& Short).
\end{cousinweb}

\begin{grammarlens}[title={What you've built}]
One more word for telling, explaining and describing at B1.
\end{grammarlens}

\subsection*{Your turn}
\begin{itemize}
\raggedright
  \item \emph{Say it:} \emph{sinistrōrsum}
  \item \emph{Say it:} \emph{sinistrōrsum}, once more
  \item \emph{Say it:} a sentence of your own with \emph{sinistrōrsum}
  \item \emph{Recall:} say the Latin for to track, then the Latin for to the right, then say \emph{sinistrōrsum} again
\end{itemize}

\begin{tcolorbox}[breakable,colback=teal!4,colframe=teal!35!black,title={Before you move on}]
What does sinistrōrsum mean? (\textquotedblleft{}To the left, towards the left\textquotedblright{}.) Say it once more.
\end{tcolorbox}
\section[oblīquus, oblīqua, oblīquum]{oblīquus, oblīqua, oblīquum — slanting, sideways, at an angle}

\label{lesson:LA-C454-obliquus}



\begin{quote}
Before the new one: say the Latin for to the right, then the Latin for to the left.
\end{quote}

\subsection*{You'll want to know: oblīquus, oblīqua, oblīquum}
\textbf{oblīquus, oblīqua, oblīquum} — \textquotedblleft{}slanting, sideways, at an angle\textquotedblright{}.

\begin{quote}
Sēmita oblīqua per collem ascendit.
\end{quote}

\begin{quote}
\emph{The path climbs across the hill at a slant.}
\end{quote}

An adjective of three endings. \emph{Oblīquō itinere}, by a sideways route. English \emph{oblique} comes from it.

\begin{cousinweb}
\emph{oblīquus} is built from \emph{ob-} and an old root meaning 'bent, slanting' (Lewis \& Short); the exact source of that root is not settled. English \emph{oblique} comes from it.
\end{cousinweb}

\begin{grammarlens}[title={What you've built}]
One more word for telling, explaining and describing at B1.
\end{grammarlens}

\subsection*{Your turn}
\begin{itemize}
\raggedright
  \item \emph{Say it:} \emph{oblīquus, oblīqua, oblīquum}
  \item \emph{Say it:} \emph{oblīquus, oblīqua, oblīquum}, once more
  \item \emph{Say it:} a sentence of your own with \emph{oblīquus, oblīqua, oblīquum}
  \item \emph{Recall:} say the Latin for to the right, then the Latin for to the left, then say \emph{oblīquus, oblīqua, oblīquum} again
\end{itemize}

\begin{tcolorbox}[breakable,colback=teal!4,colframe=teal!35!black,title={Before you move on}]
What does oblīquus, oblīqua, oblīquum mean? (\textquotedblleft{}Slanting, sideways, at an angle\textquotedblright{}.) Say it once more.
\end{tcolorbox}
\section[circumeō, circumīre]{circumeō, circumīre — to go round, to walk round; to visit in turn}

\label{lesson:LA-C454-circumeo}



\begin{quote}
Before the new one: say the Latin for to the left, then the Latin for slanting.
\end{quote}

\subsection*{You'll want to know: circumeō, circumīre}
\textbf{circumeō, circumīre} — \textquotedblleft{}to go round, to walk round; to visit in turn\textquotedblright{}.

\begin{quote}
Via mūrōs oppidī circumit.
\end{quote}

\begin{quote}
\emph{The road goes round the walls of the town.}
\end{quote}

Irregular like \emph{eō}: \emph{circumeō, circumīre, circumiī, circumitum}. Often written \emph{circueō}.

\begin{cousinweb}
\emph{circumeō} is built from \emph{circum} 'around' and \emph{eō} 'to go' (EtymDB; Lewis \& Short).
\end{cousinweb}

\begin{grammarlens}[title={What you've built}]
One more word for telling, explaining and describing at B1.
\end{grammarlens}

\subsection*{Your turn}
\begin{itemize}
\raggedright
  \item \emph{Say it:} \emph{circumeō, circumīre}
  \item \emph{Say it:} \emph{circumeō, circumīre}, once more
  \item \emph{Say it:} a sentence of your own with \emph{circumeō, circumīre}
  \item \emph{Recall:} say the Latin for to the left, then the Latin for slanting, then say \emph{circumeō, circumīre} again
\end{itemize}

\begin{tcolorbox}[breakable,colback=teal!4,colframe=teal!35!black,title={Before you move on}]
What does circumeō, circumīre mean? (\textquotedblleft{}To go round, to walk round; to visit in turn\textquotedblright{}.) Say it once more.
\end{tcolorbox}
\section[circuitus, circuitūs]{circuitus, circuitūs — a detour, a way round; a circuit}

\label{lesson:LA-C454-circuitus}



\begin{quote}
Before the new one: say the Latin for slanting, then the Latin for to go round.
\end{quote}

\subsection*{You'll want to know: circuitus, circuitūs}
\textbf{circuitus, circuitūs} — \textquotedblleft{}a detour, a way round; a circuit\textquotedblright{}.

\begin{quote}
Pōns frāctus est; itaque longō circuitū ad oppidum pervenīmus.
\end{quote}

\begin{quote}
\emph{The bridge is broken; so we reach the town by a long detour.}
\end{quote}

Fourth declension, masculine: \emph{circuitus, circuitūs m.}, from \emph{circumeō}. English \emph{circuit} comes from it.

\begin{cousinweb}
\emph{circuitus} is built on \emph{circumeō} 'to go round' (EtymDB; Lewis \& Short), with the noun ending \emph{-tus}. English \emph{circuit} comes from it.
\end{cousinweb}

\begin{grammarlens}[title={What you've built}]
One more word for telling, explaining and describing at B1.
\end{grammarlens}

\subsection*{Your turn}
\begin{itemize}
\raggedright
  \item \emph{Say it:} \emph{circuitus, circuitūs}
  \item \emph{Say it:} \emph{circuitus, circuitūs}, once more
  \item \emph{Say it:} a sentence of your own with \emph{circuitus, circuitūs}
  \item \emph{Recall:} say the Latin for slanting, then the Latin for to go round, then say \emph{circuitus, circuitūs} again
\end{itemize}

\begin{tcolorbox}[breakable,colback=teal!4,colframe=teal!35!black,title={Before you move on}]
What does circuitus, circuitūs mean? (\textquotedblleft{}A detour, a way round; a circuit\textquotedblright{}.) Say it once more.
\end{tcolorbox}
\section[(dialogue)]{Words from chapters 447–454 — a review}

\label{lesson:LA-R454-review-447-454}



\begin{quote}
Say the Latin for a bedbug and to bite.
\end{quote}

\subsection*{Your turn}
Say these aloud:

\begin{itemize}
\raggedright
  \item a bedbug, to bite, to scratch, sleepless, to lie down
  \item to snore, to make a noise, to make a din, a drinker, hoarse
  \item flat wine, neat wine, to water down, rotten, a tavern hostess
  \item a complaint, to complain, to cheat, grasping, dirt
  \item to inquire, to ask around, to ask after, to show, to make known
  \item to stray, off the road, without roads, a roundabout way, to wander about
  \item a leader, knowing, to escort, to find, to track
  \item to the right, to the left, slanting, to go round, a detour
\end{itemize}

\begin{tcolorbox}[breakable,colback=teal!4,colframe=teal!35!black,title={Before you move on}]
A bedbug? (\textbf{cīmex, cīmicis}.) To bite? (\textbf{mordeō, mordēre}.) To scratch? (\textbf{scalpō, scalpere}.) Sleepless? (\textbf{īnsomnis, īnsomne}.) To lie down? (\textbf{cubō, cubāre}.) To snore? (\textbf{stertō, stertere}.) To make a noise? (\textbf{strepō, strepere}.) To make a din? (\textbf{obstrepō, obstrepere}.) A drinker? (\textbf{pōtor, pōtōris}.) Hoarse? (\textbf{raucus, rauca, raucum}.) Flat wine? (\textbf{vappa, vappae}.) Neat wine? (\textbf{merum, merī}.) To water down? (\textbf{dīluō, dīluere}.) Rotten? (\textbf{pūtidus, pūtida, pūtidum}.) A tavern hostess? (\textbf{cōpa, cōpae}.) A complaint? (\textbf{querēla, querēlae}.) To complain? (\textbf{conqueror, conquerī}.) To cheat? (\textbf{fraudō, fraudāre}.) Grasping? (\textbf{rapāx, rapācis}.) Dirt? (\textbf{sordēs, sordis}.) To inquire? (\textbf{percontor, percontārī}.) To ask around? (\textbf{scīscitor, scīscitārī}.) To ask after? (\textbf{requīrō, requīrere}.) To show? (\textbf{mōnstrō, mōnstrāre}.) To make known? (\textbf{indicō, indicāre}.) To stray? (\textbf{aberrō, aberrāre}.) Off the road? (\textbf{dēvius, dēvia, dēvium}.) Without roads? (\textbf{invius, invia, invium}.) A roundabout way? (\textbf{ambāgēs, ambāgum}.) To wander about? (\textbf{vagor, vagārī}.) A leader? (\textbf{ductor, ductōris}.) Knowing? (\textbf{gnārus, gnāra, gnārum}.) To escort? (\textbf{prōsequor, prōsequī}.) To find? (\textbf{reperiō, reperīre}.) To track? (\textbf{vestīgō, vestīgāre}.) To the right? (\textbf{dextrōrsum}.) To the left? (\textbf{sinistrōrsum}.) Slanting? (\textbf{oblīquus, oblīqua, oblīquum}.) To go round? (\textbf{circumeō, circumīre}.) A detour? (\textbf{circuitus, circuitūs}.)
\end{tcolorbox}
